% ===================================================================
%  TABEL Nr. 9 — Die berg. — Die badplaas. — Die woud. — Die jag.
%  Meme regime que les precedents.
%
%  CETTE COLONNE SUIT L'IDO ET NON LE FRANCAIS. Le fac-simile
%  francais de ce tableau a perdu la moitie de ses intertitres et en
%  deplace deux ; l'ido a ses douze intertitres a leur rang, et le
%  neerlandais comme le tcheque les suivent. On traduit l'ido.
%
%  MEME CHOSE POUR DEUX ALINEAS QUE LE FRANCAIS ABREGE : au 5 de la
%  deuxieme scene il oublie la fuite du braconnier et le gibier reste
%  sur l'herbe (renvois 18 a 21), et au 8 il refait la phrase du
%  charbonnier. L'ido les a ; cette colonne aussi.
%
%  ET TROIS RENVOIS DE CE TABLEAU SONT CORRIGES SUR LA PAGE DE
%  LECTURE, DANS LA COLONNE FRANCAISE SEULE : elle imprime (11), (14)
%  et (19) la ou l'ido lit (41), (44) et (49). La correction vit dans
%  gravuri/korekti.json et ne touche aucune source.
%
%  DEUX PIEGES NEERLANDAIS DANS UN TABLEAU D'ARBRES : « abieti » est
%  le sapin — spar —, « pini » le pin — denneboom —, et l'afrikaans
%  ne les confond pas plus que l'ido. Et « silvo » (25) est la
%  futaie, le vieux peuplement — opgaande bos —, non le bois en
%  general, qui est « foresto ». La planche montre les deux.
% ===================================================================

%%K t09-apar-1 apar
\VUsaut{4.0mm}
\VUtitre{15.0pt}{60}{TABEL Nr. 9}
\VUsaut{3.0mm}

%%K t09-tit-1 sub
\VUcentre{12.0pt}{0}{\textit{Eerste tafereel.}}
\VUsaut{3.0mm}

%%K t09-tit-2 sub
\VUpk{11.4pt}{40}{Die berg. --- Die badplaas.}
\VUsaut{3.0mm}

%%K t09-c1-01-1 p
1. --- Ek is al agt dae in Pratz. Ek kan nie al die bewondering uitdruk wat ek gevoel het
toe ek voor die wonderlike \VUgras{bergketting} \textsuperscript{(1)} van die Pireneë
gestaan het nie, waarvan die \VUgras{kam} \textsuperscript{(2)} soos 'n reusagtige festoen
teen die blou hemel afsteek.

%%K t09-c1-02-1 p
2. --- Ek het my nie voorgestel dat die Pireneese \VUgras{toppe} \textsuperscript{(3)} so
'n groot \VUgras{hoogte} bereik nie, en ek het verstom gestaan voor die pragtige
\VUgras{gletsers} \textsuperscript{(5)} en die besneeude \VUgras{spitse}
\textsuperscript{(4)} waarop sulke verskriklike \VUgras{lawines} afrol.

%%K t09-tit-3 sub
\VUsaut{3.0mm}
\VUpk{10.2pt}{40}{Berglandskap.}
\VUsaut{3.0mm}

%%K t09-c1-03-1 p
3. --- Reeds die dag ná my aankoms het ek na die ou uitgedoofde \VUgras{vulkaan}
\textsuperscript{(6)} gegaan wat naby Pratz is. Op die dor en ontblote rande van die
\VUgras{krater} sien 'n mens nog goed die spore van die deurgang van die \VUgras{lawa}
wat etlike eeue gelede oor die \VUgras{berghang} \textsuperscript{(7)} gevloei het. Ek kon
op een van die \VUgras{hange} 'n paar brokstukke van daardie lawa vind.

%%K t09-c1-04-1 p
4. --- Pratz lê op die grens, en die \VUgras{vesting} \textsuperscript{(8)} wat die
\VUgras{bergpas} \textsuperscript{(27)} verdedig wat Frankryk van Spanje skei, herinner
heeltemal aan daardie ou kastele aan die oewer van die Ryn wat lyk of hulle van bo op
hulle rots na 'n \VUgras{prooi} loer.

%%K t09-c1-05-1 p
5. --- 'n Enorme \VUgras{arend} \textsuperscript{(9)} wat sy nes nie ver van die
\VUgras{fort} \textsuperscript{(8)} af het nie, trek dikwels my aandag. Daardie
\VUgras{roofvoël}, wie se \VUgras{snawel} sterk is, bring dikwels 'n lammetjie of 'n
bokkie na sy kleintjies wat hy met sy \VUgras{kloue} vashou. So groot is sy
\VUgras{vlerke} dat hy lank met sy swaar vrag kan sweef.

%%K t09-tit-4 sub
\VUsaut{3.0mm}
\VUpk{10.2pt}{40}{Die uitstappier.}
\VUsaut{3.0mm}

%%K t09-c1-06-1 p
6. --- Daar is die aangenaamste wandeling die uitstappie na die waterval
\textsuperscript{(10)} van «Kau». Die \VUgras{waterval} val kolkend af en verdwyn dan as
'n mooi \VUgras{spruitjie} in 'n \VUgras{sparbos} \textsuperscript{(11)} wat wes van die
stad lê. Ons het deur daardie bos opgeklim tot by die \VUgras{weerkundige observatorium}
\textsuperscript{(12)}, en van die kruin van die \VUgras{uitkykpunt} af het ons die hele
\VUgras{panorama} van die streek aanskou. Die \VUgras{landskap} is wonderlik, en die
\VUgras{natuur} lyk of sy al die skatte waaroor sy beskik aan die land uitdeel.

%%K t09-c1-07-1 p
7. --- Om af te klim het ons die \VUgras{kabelspoor} \textsuperscript{(13)} gebruik, en
ons het by 'n klein \VUgras{stasie} naby die \VUgras{ruïnes} \textsuperscript{(14)} van 'n
ou \VUgras{feodale kasteel} stilgehou wat eertyds deur die here van Pratz bewoon is.

%%K t09-c1-08-1 p
8. --- Arme kasteel! Wat dink dit seker wanneer dit onder hom die koketterige
\VUgras{badplaas} \textsuperscript{(15)} sien wat nou 'n modieuse \VUgras{termale
kuuroord} is?

%%K t09-c1-08-2 p
Die \VUgras{klimaat} is so aangenaam, die \VUgras{mineraalwater} so werksaam, dat die
\VUgras{kuur} wat in die \VUgras{sanatorium} \textsuperscript{(17)} gebruik word 'n ware
plesier is.

%%K t09-tit-5 sub
\VUsaut{3.0mm}
\VUpk{10.2pt}{40}{Aangename badplaas.}
\VUsaut{3.0mm}

%%K t09-c1-09-1 p
9. --- Origens het 'n mens alles gedoen om die verblyf aangenaam te maak. 'n Groot
\VUgras{viaduk} \textsuperscript{(16)} is gebou om 'n spoorweg moontlik te maak; die
\VUgras{park} \textsuperscript{(18)} van die \VUgras{termale inrigting}, waar
\VUgras{konserte} gehou word, is op 'n bekoorlike wyse aangelê. Etlike sierlike
«\VUgras{villa}»'s \textsuperscript{(19)} is rondom die casino \textsuperscript{(21)}
gebou, en 'n baie goed onderhoue \VUgras{grootpad} \textsuperscript{(22)} maak die
verbindings baie makliker.

%%K t09-c1-10-1 p
10. --- 'n Eenvoudige \VUgras{talud} \textsuperscript{(23)} skei die \VUgras{pad}
\textsuperscript{(20)} van die eintlike \VUgras{vallei} \textsuperscript{(25)}, in die
diepte waarvan 'n sierlike klein \VUgras{meer} \textsuperscript{(26)} lê wat 'n helder
\VUgras{spruit} \textsuperscript{(24)} voed.

%%K t09-c1-11-1 p
11. --- Aan die ander kant van die vallei word 'n \VUgras{ystermyn}
\textsuperscript{(28)} ontgin. Etlike kere het ek die \VUgras{mynwerkers} gaan sien wat
in die \VUgras{skag} afdaal, en ek het selfs die \VUgras{gang} binnegegaan waar hulle
hulle dag deurbring deur die \VUgras{erts} te ontgin wat die \VUgras{metaal} bevat.

%%K t09-c1-12-1 p
12. --- 'n Belangrike \VUgras{gietery} is naby die myn gebou om die rykdomme wat daaruit
gehaal word dadelik te benut. Die \VUgras{hoogoond} \textsuperscript{(29)}, wat met die
\VUgras{steenkool} verhit word wat hier oorvloedig is, maak dit moontlik om vinnig en
goedkoop \VUgras{gietyster} te verkry.

%%K t09-c1-13-1 p
13. --- Nie ver van die myn af grawe 'n mens 'n \VUgras{steengroef}
\textsuperscript{(30)} uit. Nóg \VUgras{graniet}, nóg \VUgras{porfier}, nóg selfs
\VUgras{sandsteen} kap die ou \VUgras{klipkapper} met sy \VUgras{houweel} los, maar
eenvoudige \VUgras{boustene} vir die bou van huise.

%%K t09-c1-14-1 p
14. --- Een van die mooiste sierade van daardie land is die \VUgras{sparre}
\textsuperscript{(31)}. Oral in die diepte van 'n \VUgras{kloof} \textsuperscript{(32)},
aan die oewer van 'n \VUgras{bergstroom} \textsuperscript{(33)}, van 'n \VUgras{afgrond}
\textsuperscript{(34)} of van 'n krans, gee hulle fyn naalde 'n groen tint.

%%K t09-tit-6 sub
\VUsaut{3.0mm}
\VUpk{10.2pt}{40}{Te voet gaan.}
\VUsaut{3.0mm}

%%K t09-c1-15-1 p
15. --- Ek maak van my vakansie gebruik om \VUgras{lank} te wandel. Die \VUgras{paaie}
\textsuperscript{(35)} wat oor die berg kronkel maak dit moontlik om, sonder om op die
afstand te let, etlike \VUgras{kilometers} en dikwels etlike \VUgras{myriameters} af te
lê.

%%K t09-c1-15-2 p
\VUgras{Mylpale} \textsuperscript{(36)} en \VUgras{wegwysers} \textsuperscript{(37)} merk
die paaie waarop 'n mens dikwels 'n jong \VUgras{bergbewoner} \textsuperscript{(38)}
teëkom met 'n \VUgras{knapsak} \textsuperscript{(40)} aan sy sy, wat 'n \VUgras{marmot}
\textsuperscript{(39)} dra, of een of ander \VUgras{sigeuner} \textsuperscript{(41)} wie
se \VUgras{pelsmus} \textsuperscript{(42)} vreemd kontrasteer met die ou geskeurde
\VUgras{kap} \textsuperscript{(43)}. Hy lei aan 'n \VUgras{ketting} 'n afgerigte
\VUgras{beer} \textsuperscript{(44)}.

%%K t09-tit-7 sub
\VUpk{10.2pt}{40}{Bergklim.}
\VUsaut{3.0mm}

%%K t09-c1-16-1 p
16. --- Byna elke dag gaan ek saam met 'n \VUgras{gids} \textsuperscript{(48)}
\VUgras{bergklim} oefen, en ek is baie trots wanneer ek met die \VUgras{sak}
\textsuperscript{(47)} op my rug en die «\VUgras{alpenstok}» in die hand, soos 'n ware
\VUgras{toeris} \textsuperscript{(46)}, die \VUgras{rotse} \textsuperscript{(45)}
bestorm.

%%K t09-c1-17-1 p
17. --- Van die kruin van 'n berg af lyk die inwoners van die vlakte nie groter as miere
nie; 'n \VUgras{trop skape} \textsuperscript{(50)} wat in 'n \VUgras{weiveld}
\textsuperscript{(49)} wei lyk soos 'n speelding vir kinders. 'n \VUgras{Denneboom}
\textsuperscript{(51)} of ook 'n \VUgras{houthuisie} \textsuperscript{(52)} lyk skaars
soos 'n vlek op die groenigheid.

%%K t09-c1-18-1 p
18. --- Tydens daardie uitstappies kom ons dikwels op die \VUgras{voetpad}
\textsuperscript{(53)} 'n \VUgras{gemsjagter} \textsuperscript{(54)} teë wat die dier wat
hy pas geskiet het op sy skouer dra.

%%K t09-c1-19-1 p
19. --- Ek het in my herbarium 'n baie merkwaardige tak \VUgras{varing}
\textsuperscript{(55)} en 'n mooi rooskleurige takkie \VUgras{heide}
\textsuperscript{(57)} gesit wat ek by die ingang van 'n klein \VUgras{grot}
\textsuperscript{(56)} tydens my laaste wandeling gepluk het.

%%K t09-tit-8 sub
\VUsaut{3.0mm}
\VUcentre{12.0pt}{0}{\textit{Tweede tafereel.}}
\VUsaut{3.0mm}

%%K t09-tit-9 sub
\VUpk{11.4pt}{40}{Die woud. --- Die jag.}
\VUsaut{3.0mm}

%%K t09-c2-01-1 p
1. --- Gister het ek 'n heerlike dag in die woud deurgebring. Die bedrywigheid wat heers
onder die diere en die \VUgras{insekte} \textsuperscript{(5)} deur wie dit bewoon word,
het my baie geïnteresseer.

%%K t09-c2-01-2 p
Die \VUgras{spinnekop} \textsuperscript{(1)} wat sy \VUgras{web} \textsuperscript{(2)}
tussen twee takke weef om die verbygaande \VUgras{vlieg} \textsuperscript{(3)} te vang,
die \VUgras{slak} \textsuperscript{(4)} wat moeisaam sy dop dra, die vlieënde hert wat na
sy prooi soek, die vlytige mier baie ywerig by die \VUgras{miershoop}
\textsuperscript{(6)}, al daardie kleintjies het geloop, gekom, gewerk met buitengewone
vlyt.

%%K t09-c2-02-1 p
2. --- 'n \VUgras{Slang} \textsuperscript{(7)} wat ek op die eerste gesig vir 'n
\VUgras{adder} aangesien het, maar wat niks anders as 'n eenvoudige \VUgras{ringslang}
was nie, het onder 'n boomstam weggekruip en het van tyd tot tyd sy fyn koppie uitgesteek
wat altyd gereed gelyk het om te byt. Voor my het 'n groot \VUgras{naakslak}
\textsuperscript{(8)} swaar gekruip nadat hy een van die smaaklike \VUgras{aarbeitjies}
\textsuperscript{(11)} verslind het wat daar oorvloedig groei.

%%K t09-c2-03-1 p
3. --- Die grond was met \VUgras{bosblomme} bekleed; \VUgras{sleutelblomme}
\textsuperscript{(9)}, \VUgras{maagdepalms} \textsuperscript{(10)} en baie ander plante
wat ek nie geken het nie, het tussen \VUgras{grasbossies} \textsuperscript{(12)} geblom.

%%K t09-c2-04-1 p
4. --- In daardie streek is die \VUgras{truffel} byna so algemeen soos die
\VUgras{sampioen} \textsuperscript{(14)}, en 'n \VUgras{varkie} \textsuperscript{(13)}
wat aan 'n boswagter behoort, het gulsig gesoek of hy 'n paar daarvan sou kon vind.

%%K t09-tit-10 sub
\VUsaut{3.0mm}
\VUpk{10.2pt}{40}{Wildstropery.}
\VUsaut{3.0mm}

%%K t09-c2-05-1 p
5. --- Ek het die \VUgras{einde} bygewoon van 'n tafereel wat my baie vermaak het. Met
behulp van die reuk van sy \VUgras{basset} \textsuperscript{(15)} het die
\VUgras{boswagter} \textsuperscript{(16)} pas 'n \VUgras{wildstroper}
\textsuperscript{(22)} oorval op die oomblik toe dié hom gereedgemaak het om die
\VUgras{wild} \textsuperscript{(17)} saam te neem wat hy doodgemaak het. Omdat hy eerder
sy buit wou prysgee as om hom te laat aankeer, het die wildstroper gevlug, en 'n
\VUgras{haas} \textsuperscript{(18)}, 'n \VUgras{hinde} \textsuperscript{(19)}, 'n
\VUgras{ree} \textsuperscript{(20)} en 'n \VUgras{wildevark} \textsuperscript{(21)} het op
die gras gelê tot vreugde van die boswagter.

%%K t09-c2-06-1 p
6. --- Die verskeidenheid van die bome wat die woud bevat is verstommend. Die
\VUgras{beuke} \textsuperscript{(23)} en die \VUgras{berke} \textsuperscript{(26)}, die
\VUgras{dennebome}, wat die \VUgras{harpuis} \textsuperscript{(27)} lewer, die
\VUgras{eike} \textsuperscript{(29)} en ook baie ander meng hulle takke.

%%K t09-c2-06-2 p
Die \VUgras{klimop} \textsuperscript{(24)}, wat aan die stam van die hoë bome vasklou,
kleur die \VUgras{opgaande bos} \textsuperscript{(25)} met blinkende groenigheid, wat hom
aangenaam verenig met die helderder en bleekgroen tint van die \VUgras{mos}
\textsuperscript{(31)} waarin die \VUgras{groenspeg} \textsuperscript{(30)} insekte soek.

%%K t09-tit-11 sub
\VUsaut{3.0mm}
\VUpk{10.2pt}{40}{Woudlandskap.}
\VUsaut{3.0mm}

%%K t09-c2-07-1 p
7. --- Die ou eike het my bekoor. Hulle groot verwronge \VUgras{wortels}
\textsuperscript{(32)}, wat half uit die grond gekom het, gee hulle 'n pragtige voorkoms
van krag en fleur, en ek vra myself af hoe die \VUgras{eienaar} \textsuperscript{(35)}
dit oor sy hart kan kry om hulle te laat afkap en aan die \VUgras{saag}
\textsuperscript{(34)} van die \VUgras{houtkapper} \textsuperscript{(33)} oor te lewer.
Die arme \VUgras{stamme} \textsuperscript{(36)}, gestroop van hulle \VUgras{bas}
\textsuperscript{(37)} of gekloof deur die \VUgras{byl} \textsuperscript{(38)} van die
\VUgras{dagloner} \textsuperscript{(39)}, bedroef my.

%%K t09-c2-08-1 p
8. --- Naby het 'n ou \VUgras{houtskoolbrander} \textsuperscript{(41)} met sy
\VUgras{graaf 'n hoop} \textsuperscript{(42)} houtskool voorberei, en 'n ou vrou het 'n
\VUgras{bondel} \textsuperscript{(40)} \VUgras{dooie hout} van takkies van die afgekapte
bome saamgedra.

%%K t09-tit-12 sub
\VUsaut{3.0mm}
\VUpk{10.2pt}{40}{Jag.}
\VUsaut{3.0mm}

%%K t09-c2-09-1 p
9. --- Ek wou 'n paar oomblikke in die \VUgras{boswagtershuis} \textsuperscript{(46)}
gaan rus, maar toe ek naby die \VUgras{dammetjie} \textsuperscript{(43)} gekom het waarop
'n mooi \VUgras{swaan} \textsuperscript{(45)} swem, het ek opgemerk dat 'n onlangse
\VUgras{oorstroming} \textsuperscript{(44)} die pad in 'n ware \VUgras{moeras} verander
het. Ek moes omdraai, en toe ek verby die \VUgras{seder} \textsuperscript{(49)}, die
\VUgras{populiere} \textsuperscript{(48)} en die \VUgras{olm} \textsuperscript{(47)}
gegaan het wat aan die rand van die woud staan, het ek uit 'n \VUgras{kreupelbos}
\textsuperscript{(54)} 'n pragtige \VUgras{takbok} \textsuperscript{(50)} sien uitkom,
agtervolg deur 'n hardnekkige \VUgras{meute} \textsuperscript{(51)} wat ook deur die
\VUgras{horing} \textsuperscript{(53)} van die \VUgras{jagters te perd}
\textsuperscript{(52)} opgesweep is. Die arme dier was heeltemal uitgeput. Dit het
sterwend 'n paar treë van my af neergeval.

%%K t09-c2-10-1 p
10. --- Die seun van die boswagter, wat deur die lawaai van die \VUgras{jag te perd}
gelok is, het uit die huis gekom en ons twee het na die woud teruggegaan. Hy het 'n
\VUgras{ekster} \textsuperscript{(56)} op 'n tak van 'n ou eik gesien. Hy het gemeen dat
haar nes seker in die \VUgras{loof} \textsuperscript{(58)} verborge is, en hy wou die nes
vang.

%%K t09-c2-10-2 p
Rats soos 'n \VUgras{eekhoring} \textsuperscript{(61)} het hy van \VUgras{tak}
\textsuperscript{(55)} tot tak geklim, maar hy het niks gevind nie en het net daarin
geslaag om 'n ou \VUgras{kraai} \textsuperscript{(60)} te laat opvlieg wat op 'n
\VUgras{takkie} \textsuperscript{(57)} gesit het.
\VUsaut{4.0mm}
\VUfilet{22mm}
